\apendice{Plan de Proyecto Software}

\section{Introducción}

Este anexo describe la planificación del proyecto, el seguimiento realizado con Zube y el análisis de viabilidad del sistema.

\section{Planificación temporal}

El desarrollo del proyecto se ha organizado siguiendo un enfoque incremental basado en sprints. El seguimiento se ha realizado en Zube, utilizando historias de usuario, tareas técnicas y cambios asociados a commits del repositorio. El tablero del proyecto se encuentra en:

\begin{center}
\url{https://zube.io/tfg-azure-servicio-empresarial/tfg-servicio-empresarial-seguro/w/workspace-1/kanban}
\end{center}

Las principales fases del proyecto han sido:

\begin{itemize}

\item Análisis del problema y definición de requisitos.

\item Diseño de la arquitectura del sistema.

\item Implementación del prototipo.

\item Despliegue del sistema en la plataforma cloud.

\item Validación y pruebas del sistema.

\end{itemize}

\subsection{Iteraciones realizadas}

La planificación se ha dividido en cinco iteraciones, superando el mínimo de tres sprints indicado en la guía docente. Cada iteración generó un incremento funcional o documental verificable:

\begin{table}[h]
\centering
\small
\begin{tabularx}{\textwidth}{l X X}
\textbf{Sprint} & \textbf{Objetivo} & \textbf{Resultado} \\
\hline
Sprint 1 (25/02/2026--10/03/2026) & Preparación y base. & Repositorio, documentación inicial, backlog y primeras decisiones de arquitectura. \\
Sprint 2 (11/03/2026--24/03/2026) & Implementación de módulos y dashboard. & Clasificación ligera, portal inicial, persistencia y primeras pruebas unitarias. \\
Sprint 3 (25/03/2026--21/04/2026) & Incidencias y despliegue. & API de solicitudes, validación, seguridad básica y despliegue en Azure. \\
Sprint 4 (22/04/2026--21/05/2026) & Consolidación y documentación técnica. & App Service, Storage, Key Vault, Managed Identity, evidencias y anexos. \\
Sprint 5 (22/05/2026--08/07/2026) & Finalización y defensa. & Memoria, anexos, calidad SonarCloud, vídeos y entregables finales. \\
\end{tabularx}
\caption{Resumen de iteraciones del proyecto}
\end{table}

\subsection{Seguimiento en Zube}

Zube se ha utilizado como tablero Kanban para registrar historias de usuario, tareas técnicas y evolución por sprint. La tabla \ref{tabla:seguimiento_zube} resume el estado observable del seguimiento realizado.

\begin{table}[h]
\centering
\small
\begin{tabularx}{\textwidth}{l l l X}
\textbf{Sprint} & \textbf{Inicio} & \textbf{Cierre/estado} & \textbf{Evidencia de seguimiento} \\
\hline
Sprint 1 & 25/02/2026 & Cerrado 15/03/2026 & 6 tarjetas cerradas; preparación, repositorio, arquitectura inicial y SonarCloud previsto. \\
Sprint 2 & 11/03/2026 & Cerrado 17/04/2026 & 6 tarjetas cerradas; módulos, dashboard, IA ligera y pruebas unitarias iniciales. \\
Sprint 3 & 25/03/2026 & Cerrado 20/05/2026 & Integración de incidencias, seguridad, despliegue y validación de endpoints. \\
Sprint 4 & 22/04/2026 & Cerrado 22/05/2026 & 5 tarjetas cerradas; consolidación técnica, Azure, documentación y evidencias. \\
Sprint 5 & 22/05/2026 & Abierto hasta 08/07/2026 & Tareas finales de memoria, anexos, vídeos, calidad, release y entrega UBUVirtual. \\
\end{tabularx}
\caption{Seguimiento de sprints registrado en Zube}
\label{tabla:seguimiento_zube}
\end{table}

Las capturas asociadas al seguimiento se conservan en la carpeta \texttt{docs/sprints/}. El Sprint 5 permanece abierto porque agrupa tareas finales de defensa y entrega oficial que se cierran una vez generados los vídeos, el PDF de enlaces y la release final.

\section{Estudio de viabilidad}

\subsection{Viabilidad económica}

El proyecto se ha desarrollado utilizando principalmente herramientas disponibles en versiones gratuitas, de código abierto o académicas, como GitHub, Zube y Microsoft Azure for Students. SonarCloud se utiliza como herramienta de análisis de calidad del código y su panel se aporta como evidencia de revisión externa.

La mayor parte del coste corresponde al esfuerzo de desarrollo. Para estimar el coste de personal se ha considerado una dedicación aproximada de 300 horas y un coste de referencia de 18 euros/hora para un perfil junior de desarrollo cloud. Sobre este coste se añade una estimación del 32\% de costes empresariales asociados a Seguridad Social y otras aportaciones obligatorias. Esta cifra no representa una factura real, sino una estimación académica del valor del trabajo realizado.

\begin{table}[h]
\centering
\small
\begin{tabularx}{\textwidth}{X r r r}
\textbf{Concepto} & \textbf{Cantidad} & \textbf{Coste unitario} & \textbf{Coste estimado} \\
\hline
Análisis y planificación & 45 h & 18 EUR/h & 810 EUR \\
Diseño de arquitectura & 45 h & 18 EUR/h & 810 EUR \\
Implementación backend y portal & 90 h & 18 EUR/h & 1.620 EUR \\
Infraestructura Azure y despliegue & 60 h & 18 EUR/h & 1.080 EUR \\
Pruebas, documentación y memoria & 60 h & 18 EUR/h & 1.080 EUR \\
Costes empresariales de personal & 32\% & 5.400 EUR & 1.728 EUR \\
Servicios Azure for Students & 1 suscripción & 0 EUR & 0 EUR \\
GitHub y Zube & 1 cuenta académica & 0 EUR & 0 EUR \\
\hline
\textbf{Total estimado} & \textbf{300 h} & & \textbf{7.128 EUR} \\
\end{tabularx}
\caption{Estimación económica del proyecto}
\end{table}

En un entorno empresarial real también deberían considerarse costes de operación continuada, soporte, monitorización, almacenamiento y transferencia de datos. En este TFG se han mantenido los recursos dentro de un entorno de desarrollo de bajo consumo, por lo que el coste de infraestructura se ha reducido al mínimo. El equipo informático y el software de escritorio no se imputan por su coste completo, ya que conservan valor tras el proyecto y pueden reutilizarse en otros trabajos; en una contabilidad empresarial se imputaría únicamente su amortización proporcional al periodo de uso.

\subsection{Viabilidad legal}

Durante el desarrollo del proyecto se han utilizado herramientas de software bajo licencias abiertas o gratuitas para uso académico. Además, no se han utilizado datos personales reales, por lo que no se han generado problemas relacionados con la protección de datos.

\begin{table}[h]
\centering
\small
\begin{tabularx}{\textwidth}{l l X}
\textbf{Componente} & \textbf{Licencia / modalidad} & \textbf{Uso en el proyecto} \\
\hline
Python & PSF License & Lenguaje backend \\
Flask & BSD-3-Clause & API REST y portal web \\
Gunicorn & MIT & Servidor WSGI en App Service \\
Azure SDK for Python & MIT & Integración con Storage y Key Vault \\
Terraform & Business Source License & Infraestructura como código \\
Azure CLI & Licencia Microsoft & Automatización del despliegue desde PowerShell \\
SonarCloud & Servicio SonarSource & Análisis de calidad del código \\
Microsoft Azure & Servicio cloud & App Service, Storage, Key Vault y Managed Identity \\
\end{tabularx}
\caption{Licencias y servicios utilizados}
\end{table}

Las licencias empleadas son compatibles con un proyecto académico y permiten la distribución del código fuente del prototipo. El repositorio incluye un archivo \texttt{LICENSE} con licencia MIT para el código propio desarrollado durante el TFG. Las credenciales y secretos no se almacenan en el repositorio, sino en variables de entorno y Azure Key Vault.
